\documentclass[11pt]{article}
\usepackage{tutorial}
\begin{document}
\tutorialtitle{3}{Gradient Estimators and Variance Reduction}{Score functions, control variates, importance sampling, Rao--Blackwell, and unbiased sparse credit}

\section*{Why this tutorial}
Hard decisions (which route wins, which memory slot is written) have no ordinary derivative, and exact credit is often too
expensive to compute everywhere. Learning then relies on \emph{estimators} of the gradient. The two questions are always:
is it \emph{unbiased} (right on average), and what is its \emph{variance} (how noisy)? This tutorial covers the standard toolkit
and the specific construction our sparse-credit learners use: a cheap prediction of the credit, corrected by random audits so
that it stays exactly unbiased.

\section{The score-function estimator}
For a distribution $p_\theta(z)$ and a function $f$,
\[
\nabla_\theta\,\E_{z\sim p_\theta}[f(z)]=\E_{z\sim p_\theta}\big[f(z)\,\nabla_\theta\log p_\theta(z)\big].
\]
\emph{Proof:} $\nabla\!\int p f=\int f\nabla p=\int f\,p\,\nabla\log p$. One sample $z$ gives the unbiased estimate
$f(z)\nabla\log p_\theta(z)$ (REINFORCE). It needs no derivative of $f$, so it works for discrete choices, but its variance can be large.

\paragraph{Baselines (control variates).} Because $\E[\nabla\log p]=0$, subtracting any constant $b$ keeps it unbiased:
$(f(z)-b)\nabla\log p$. For a scalar parameter the variance-minimizing constant is
$b^\star=\E[f\,(\nabla\log p)^2]/\E[(\nabla\log p)^2]$.
\begin{keyidea}
A control variate is anything with known expectation that is correlated with your estimate. Subtracting it costs no bias and
removes the shared noise. Better predictions of the quantity give lower variance; a perfect prediction gives zero variance.
\end{keyidea}

\section{Importance sampling and self-normalization}
To estimate $\E_p[g]$ with samples from a proposal $q$: $\E_p[g]=\E_q[w\,g]$ with $w=p/q$, unbiased if $q>0$ wherever $pg\ne0$.
When $p$ is known only up to a constant (a posterior), one uses \emph{self-normalized} weights $\tilde w_s=w_s/\sum_r w_r$.
This is biased at finite sample size $S$ (bias $O(1/S)$ under moment conditions) but consistent. With $S=1$ the normalized
weight is always 1: it carries no information about whether the proposal was good.

\section{Rao--Blackwellization}
If $\hat g$ is unbiased and $Y$ is any available information, then $\E[\hat g\mid Y]$ is also unbiased and
\[
\Var(\E[\hat g\mid Y])\le\Var(\hat g)
\]
(law of total variance: $\Var\hat g=\E\Var(\hat g\mid Y)+\Var\E[\hat g\mid Y]$). Analytically averaging over anything you can
average over never hurts. Example: at the output of a race, use the exact posterior responsibilities $\rho_i$ rather than one
sampled winner.

\section{Latent causes: Fisher's identity and the missing information}
For a model with a latent cause $k$ (which route, which clock), the marginal likelihood gradient is
\[
\nabla\log p(y)=\E_{k\sim p(k\mid y)}\big[\nabla\log p(y,k)\big]\qquad\text{(Fisher's identity).}
\]
So a single \emph{posterior} sample $k$ gives an unbiased estimate. Its covariance is the \emph{missing-information matrix}
$\Cov_{k\mid y}[\nabla\log p(y,k)]$. If every complete-data score has norm at most $G$ and $\rho_k=p(k\mid y)$,
\[
\tr\Cov\le 2G^2\Big(1-\sum_k\rho_k^2\Big)=2G^2\big(1-e^{-H_2(\rho)}\big),
\]
where $H_2$ is the collision (R\'enyi-2) entropy. \emph{Proof sketch:} the variance of a discrete distribution is half its mean
squared pairwise distance; each squared distance is at most $(2G)^2$ and pairs with $k\ne l$ have total weight $1-\sum\rho_k^2$.
\emph{Reading:} sampled-cause credit is nearly exact when the posterior is sharp.

\section{Predicted credit plus sparse audits (unbiased sparse credit)}
Let $g_b$ be the exact credit for block $b$ (a parameter block or a state cotangent) and $h_b$ a prediction of it, made before the
audit from information already available. Audit block $b$ independently with probability $p_b>0$ (indicator $I_b$):
\[
\hat g_b=h_b+\frac{I_b}{p_b}\,(g_b-h_b).
\]
\begin{proposition}
$\E[\hat g_b]=g_b$ and $\E\|\hat g_b-g_b\|^2=\big(\tfrac1{p_b}-1\big)\|g_b-h_b\|^2$.
\end{proposition}
\begin{proof}
$\E[I_b/p_b]=1$ gives the mean. The error is $(I_b/p_b-1)(g_b-h_b)$ and $\E(I_b/p_b-1)^2=1/p_b-1$.
\end{proof}
The prediction affects only the variance, never the mean. A good predictor makes audits rarely necessary.

\paragraph{Where to audit (Neyman allocation).} With audit costs $c_b$ and a variance budget $V$, minimize $\sum_b c_bp_b$ subject
to $\sum_b\|r_b\|^2/p_b\le V$, where $r_b=g_b-h_b$ is the \emph{surprise}. By Cauchy--Schwarz,
\[
\sum_b c_bp_b\;\ge\;\frac{\big(\sum_b\|r_b\|\sqrt{c_b}\big)^2}{V},\qquad\text{with equality at }p_b\propto\frac{\|r_b\|}{\sqrt{c_b}}
\]
(capped at $p_b\le1$). Credit cost therefore scales with the \emph{square} of the surprise: heavy while the predictor is poor,
cheap once surprises are small, and heavy again if the domain shifts.

\begin{remark}
The allocation must use only information available \emph{before} the audit draw (a predicted surprise), and the predictor must
be frozen for the current draw. Re-fitting the predictor on the audited value and then recomputing the current baseline breaks
the conditioning argument and can bias the estimate.
\end{remark}

\section*{Exercises}
\begin{exercise} Show that $\E_{p_\theta}[\nabla\log p_\theta]=0$.\end{exercise}
\begin{solution} $\int p\nabla\log p=\int\nabla p=\nabla\!\int p=\nabla1=0$.\end{solution}
\begin{exercise} In the audit estimator, what happens with $p_b=1$? With $h_b=g_b$?\end{exercise}
\begin{solution} $p_b=1$: $\hat g_b=g_b$ exactly (always audit). $h_b=g_b$: zero variance for any $p_b$ (no surprise).\end{solution}
\begin{exercise} Two causes with posterior $(0.9,0.1)$ and scores bounded by $G=1$. Bound the variance of single-sample credit.\end{exercise}
\begin{solution} $1-\sum\rho^2=1-0.82=0.18$, so $\tr\Cov\le0.36$. For $(0.5,0.5)$ it would be $1.0$.\end{solution}
\begin{exercise} Derive the Neyman allocation with Lagrange multipliers.\end{exercise}
\begin{solution} Stationarity of $\sum c_bp_b+\mu(\sum\|r_b\|^2/p_b-V)$: $c_b=\mu\|r_b\|^2/p_b^2$, so $p_b=\sqrt\mu\,\|r_b\|/\sqrt{c_b}$.\end{solution}

\begin{inourmodels}
REINFORCE, straight-through and closed-loop baselines in the hard-routing credit diagnostics (note 160 \S\S3--9); the collision-entropy
bound (note 164 Theorem F, Lean-checked); forward-predicted credit with unbiased residual audits (notes 162--163); surprise-squared
credit cost (note 165 \S10, Lean-checked); Rao--Blackwell conditioning on forward information (note 164 Result 1).
\end{inourmodels}
\end{document}
